% Source: https://artofproblemsolving.com/community/c3h1661672_points_near
% Stable identity: aops-topic-1661672-q01
\textbf{(a)} A point of the curve is $(x,x^2)$; its squared distance from $(3,0)$ is
\[
g(x)=(x-3)^2+x^4 .
\]
Minimising the distance is the same as minimising $g$. Now
\[
g'(x)=2(x-3)+4x^3=2\left(2x^3+x-3\right)=2(x-1)\left(2x^2+2x+3\right).
\]
The quadratic factor $2x^2+2x+3$ has discriminant $4-24<0$, so it is always positive. Hence $g'(x)<0$ for $x<1$ and $g'(x)>0$ for $x>1$: $g$ decreases and then increases, and its minimum is at $x=1$. The closest point is
\[
\boxed{P=(1,1)},
\]
at distance $\sqrt{g(1)}=\sqrt5$ from $(3,0)$.

\textbf{(b)} At the minimising point $x=a$ we have $g'(a)=0$, that is
\[
2(a-3)+2\cdot b\cdot 2a=0\qquad(\text{using }b=a^2),
\]
which says $(a-3,\,b)\cdot(1,\,2a)=0$. The vector $(a-3,b)$ points along the line from $(3,0)$ to $(a,b)$, and $(1,2a)$ points along the tangent at $(a,b)$, since the tangent has slope $\dfrac{dy}{dx}=2a$. Their dot product is $0$, so the line through $(3,0)$ and $P$ is perpendicular to the tangent at $P$.

For $P=(1,1)$ concretely: the tangent has slope $2$ and the line $PQ$ to $Q=(3,0)$ has slope $\dfrac{0-1}{3-1}=-\dfrac{1}{2}$, and $2\cdot\left(-\dfrac{1}{2}\right)=-1$.

\begin{asy}
usepackage("amsmath");
size(6cm);
defaultpen(fontsize(10pt));
import graph;
real f(real x){return x^2;}
draw((-1.6,0)--(3.6,0),gray,Arrow(TeXHead));draw((0,-0.6)--(0,3.8),gray,Arrow(TeXHead));
label("$x$",(3.6,0),E,gray);label("$y$",(0,3.8),N,gray);
draw(graph(f,-1.4,1.9),linewidth(0.9));
pair P=(1,1),Q=(3,0);
draw((0.3,-0.4)--(2.15,3.3),dashed);
draw(P--Q,red+linewidth(0.9));
pair u=unit((1,2)),w=unit(Q-P);real s=0.2;
draw(P+s*u--P+s*u+s*w--P+s*w);
dot(P);dot(Q);
label("$P(1,1)$",P,W);label("$(3,0)$",Q,S);
label("tangent",(2.15,3.3),E);
label("\mbox{}",(min(currentpicture,user=true)+max(currentpicture,user=true))/2,black); // keeps the text after the figure black when LaTeX typesets the labels inline
\end{asy}
